(1) 中心展开 \(k(\omega_0+\Omega)=k_0+k_1\Omega+k_2\Omega^2/2+O(\Omega^3)\)，其中 \(k_0=n_0\omega_0/c\)、\(k_1=[n_0+\omega_0n_0']/c\)、\(k_2=[2n_0'+\omega_0n_0'']/c\)，撇号均对角频率求导。

(2) 穿过长度 \(L\) 的频谱获得相位 \(k(\omega)L\)，其线性相位使包络延迟 \(\tau_g=k_1L=L(n_0+\omega_0n_0')/c\)。相对同长度真空路径的额外时延为 \(\boxed{\Delta\tau_g=L(n_g-1)/c}\)，其中群折射率 \(n_g=n_0+\omega_0n_0'\)；相速是 \(c/n_0\)，两者不可混淆。

(3) 去掉纯相位和群时延后，剩余二次色散系数为 \(\beta_2=k_2\)。Gaussian频谱乘单位模相位 \(e^{i\beta_2L\Omega^2/2}\)，逆变换后强度包络保持Gaussian但时间宽度变成 \(\boxed{\tau_L=\tau_0\sqrt{1+(\beta_2L/\tau_0^2)^2}}\)。脉冲总积分能量因Parseval严格守恒。

(4) 定义色散长度 \(L_D=\tau_0^2/|\beta_2|\)，当 \(L\ll L_D\) 时展宽很小；当 \(L\gg L_D\) 时显著展宽。若输入已带啁啾，则二次色散可先抵消初始相位的频率—时间相关而压缩，选择正确符号才能实现。高阶色散与吸收被本题忽略。
